% optimizers_explained_sorted.tex
\documentclass[12pt]{article}

\usepackage[utf8]{inputenc}
\usepackage[T1]{fontenc}
\usepackage{lmodern}
\usepackage{geometry}
\geometry{a4paper, margin=2.5cm}

\usepackage{amsmath, amsfonts, amssymb}
\usepackage{bm}
\usepackage{xcolor}
\usepackage{listings}

% --------------------------------------------------------------------
\lstset{
  language=Python,
  basicstyle=\ttfamily\small,
  frame=single,
  breaklines=true,
  tabsize=2,
  keywordstyle=\color{blue},
  commentstyle=\color{gray},
  stringstyle=\color{orange}
}
% --------------------------------------------------------------------

\title{Optimization Algorithms in Neural--Network Training\\(Ordered by Increasing Complexity)}
\author{}
\date{}

\begin{document}
\maketitle

The optimizers are presented from the \emph{least} to the \emph{most}
sophisticated according to the amount of state they keep and the number of
computations performed per iteration.

%---------------------------------------------------------------------
\section*{Common Notation}
%---------------------------------------------------------------------
\begin{itemize}
  \item[$\bm x_t$] Parameter (weight) vector at iteration~$t$.
  \item[$f(\bm x)$] Objective (loss) function to be minimised.
  \item[$\nabla_{\!\bm x}f(\bm x)$] Gradient of $f$ with respect to~$\bm x$.
  \item[$\bm g_t$] Shorthand for $\nabla_{\!\bm x}f(\bm x_t)$.
  \item[$\alpha,\eta$] Base learning rate (step size).
  \item[$\gamma$] Decay factor for the second--moment accumulator (RMSProp, backtracking).
  \item[$\beta$] Momentum coefficient.
  \item[$\beta_1,\beta_2$] Exponential decay rates for first and second moments (Adam).
  \item[$\varepsilon$] Small constant to prevent division by~$0$.
  \item[$\bm v_t$] Exponential moving average of past gradients (``momentum'').
  \item[$\bm s_t$] Exponential moving average of squared gradients (``variance'').
  \item[$\odot$] Hadamard (element--wise) product; $\bm g_t^{\odot 2}$ squares each coordinate.
\end{itemize}

%=====================================================================
\section{Vanilla Gradient Descent (Fixed Step)}
%---------------------------------------------------------------------
\subsection*{Update rule}
\begin{align}
\bm g_t      &=\nabla_{\!\bm x}f(\bm x_t), &
\bm d_t &= -\bm g_t \\
\bm x_{t+1}  &=\bm x_t+\alpha\,\bm d_t
\end{align}
with constant step size $\alpha$.

\subsection*{Rationale}
A direct step opposite to the gradient.  Simplicity is its strength, but the
fixed $\alpha$ must be carefully tuned: too small slows convergence, too large
causes divergence.

\subsection*{Python implementation}
\begin{lstlisting}
def gradient_descent(x0, f, grad, step_size,
                     niter=500, tol=1e-6):
    x = x0
    for t in range(niter):
        g = grad(x)
        if np.linalg.norm(g) <= tol:
            break
        x = x - step_size * g
    return x
\end{lstlisting}

%=====================================================================
\section{Gradient Descent with Adaptive Step\\(Backtracking Line Search)}
%---------------------------------------------------------------------
\subsection*{Update rule}
\begin{enumerate}
  \item $\bm d_t = -\,\nabla_{\!\bm x}f(\bm x_t)$
  \item Find $\alpha$ by backtracking:
        \[
        \text{while } f(\bm x_t+\alpha\bm d_t) > f(\bm x_t)
        \text{ set } \alpha \leftarrow \gamma\alpha
        \]
  \item $\bm x_{t+1} = \bm x_t + \alpha \bm d_t$
\end{enumerate}

\subsection*{Rationale}
Automatically shrinks the step until a decrease in the loss is observed,
removing most of the learning--rate tuning burden while keeping the method
conceptually simple.

\subsection*{Python implementation}
\begin{lstlisting}
def gradient_descent_adaptive_step(x0, f, grad, step_size,
                                   niter=500, tol=1e-6, gamma=0.98):
    x = x0
    for t in range(niter):
        g = grad(x)
        if np.linalg.norm(g) <= tol:
            break
        d = -g
        alpha = step_size
        while f(x + alpha * d) > f(x):
            alpha *= gamma
        x = x + alpha * d
    return x
\end{lstlisting}

%=====================================================================
\section{Gradient Descent with Momentum}
%---------------------------------------------------------------------
\subsection*{Update rule}
\begin{align}
\bm g_t  &= \nabla_{\!\bm x}f(\bm x_t) \\
\bm v_t  &= \beta\,\bm v_{t-1} + \bm g_t \\
\bm x_{t+1} &= \bm x_t - \eta\,\bm v_t
\end{align}

\subsection*{Rationale}
The velocity $\bm v_t$ accumulates a decaying history of past gradients,
acting like physical momentum.  This accelerates learning along consistent
directions and dampens oscillations across steep ravines.

\subsection*{Python implementation}
\begin{lstlisting}
def gradient_descent_momentum(x0, f, grad, niter=500, tol=1e-6,
                              eta=0.9, beta=0.98):
    x = x0
    v = np.zeros_like(x0)
    for t in range(niter):
        g = grad(x)
        if np.linalg.norm(g) <= tol:
            break
        v = beta * v + g
        x = x - eta * v
    return x
\end{lstlisting}

%=====================================================================
\section{RMSProp}
%---------------------------------------------------------------------
\subsection*{Update rule}
\begin{align}
\bm g_t &= \nabla_{\!\bm x}f(\bm x_t) \\
\bm s_t &= \gamma\,\bm s_{t-1} + (1-\gamma)\,\bm g_t^{\odot 2} \\
\bm x_{t+1} &= \bm x_t - \eta\,\dfrac{\bm g_t}{\sqrt{\bm s_t}+\varepsilon}
\end{align}

\subsection*{Rationale}
Keeps a moving average of squared gradients so each coordinate receives its
own adaptive learning rate: large variances produce small effective steps and
\emph{vice versa}.  Particularly effective in ``ravine'' landscapes with very
different curvatures along different axes.

\subsection*{Python implementation}
\begin{lstlisting}
def rmsprop(x0, f, grad, niter=500, tol=1e-6,
            eta=0.9, gamma=0.98):
    epsilon = 1e-8
    x = x0
    s = np.zeros_like(x0)
    for t in range(niter):
        g = grad(x)
        if np.linalg.norm(g) <= tol:
            break
        s = gamma * s + (1 - gamma) * g**2
        x = x - eta * g / (np.sqrt(s) + epsilon)
    return x
\end{lstlisting}

%=====================================================================
\section{Adam}
%---------------------------------------------------------------------
\subsection*{Update rule}
\begin{align}
\bm g_t      &= \nabla_{\!\bm x}f(\bm x_t) \\
\bm v_t      &= \beta_1 \bm v_{t-1} + (1-\beta_1)\,\bm g_t \\
\bm s_t      &= \beta_2 \bm s_{t-1} + (1-\beta_2)\,\bm g_t^{\odot 2} \\
\hat{\bm v}_t&= \bm v_t / (1-\beta_1^{t}),\quad
\hat{\bm s}_t= \bm s_t / (1-\beta_2^{t}) \\[4pt]
\bm x_{t+1}  &= \bm x_t -
               \eta\,\dfrac{\hat{\bm v}_t}{\sqrt{\hat{\bm s}_t}+\varepsilon}
\end{align}

\subsection*{Rationale}
Combines momentum (first moment) and per--coordinate variance scaling (second
moment) while correcting early--iteration bias.  Adam therefore adapts both
direction and magnitude of updates, offering fast convergence with minimal
hyperparameter tuning in practice.

\subsection*{Python implementation}
\begin{lstlisting}
def adam(x0, f, grad, niter=500, tol=1e-6,
         eta=0.9, beta1=0.9, beta2=0.99):
    epsilon = 1e-8
    x = x0
    v = np.zeros_like(x0)
    s = np.zeros_like(x0)
    t = 1
    while (t < niter) and (np.linalg.norm(grad(x)) > tol):
        g  = grad(x)
        v  = beta1 * v + (1 - beta1) * g
        s  = beta2 * s + (1 - beta2) * g**2
        vt = v / (1 - beta1**t)
        st = s / (1 - beta2**t)
        x  = x - eta * vt / (np.sqrt(st) + epsilon)
        t += 1
    return x
\end{lstlisting}

\end{document}
